\input format.tex

\usepackage{graphicx}
\graphicspath{{cores/}}

\usepackage{environ}
\usepackage{colortbl,array,booktabs}
\usepackage{tabularx}

\colorlet{TablaBordeSuperior}{topcolor}
\colorlet{TablaBordeInferior}{topcolor}
\colorlet{TablaCentroSuperior}{blue!1}
\colorlet{TablaCentroInferior}{blue!20}
\colorlet{FuenteCabeceraTabla}{white}

\newcolumntype{M}[1]{>{\centering\arraybackslash}m{#1}}
%\newcommand{\tabularxcolumn}[1]{>\arraybackslash}m{#1}}

\tcbset{rtab/.style={
freelance,
frame code={
 \path[top color=topcolor,bottom color=topcolor]
   ([yshift=-#1*(\baselineskip+2pt)]interior.north west) --
   ([yshift=-#1*(\baselineskip+2pt)]interior.north east) {[sharp corners]--
    ([yshift=3pt]interior.north east) --
    ([yshift=3pt]interior.north west)} -- cycle;

  },
interior code={},
 }
}

\newcommand\fuentecabecera[1]{\textcolor{black}{\textbf{#1}}}

\begin{document}

\vspace*{4mm}
%% 各章节
\setlength{\arrayrulewidth}{.2pt}
\fontsize{8.8pt}{11pt}\selectfont
\color{black100}

\vspace*{3mm}

\noindent{\fontsize{15pt}{11pt}\selectfont{ 三、主要有益菌检测结果}}

%\vspace*{-4.25mm}
\fontsize{9pt}{11pt}\selectfont
\arrayrulecolor{gray2}
\begin{longtable}{|m{4.8cm}<{\centering}|m{1.3cm}<{\centering}|m{1.3cm}<{\centering}|m{1.3cm}<{\centering}|m{4.9cm}<{\centering}|}
%\begin{longtable}{|m{4cm}<{\centering}|m{1.5cm}<{\centering}|m{1.5cm}<{\centering}|m{1.5cm}<{\centering}|m{4.9cm}<{\centering}|}
\hline
\parbox[c]{\hsize}{\vskip10pt\centering\bf {检测内容} \vskip10pt} & \bf{检测结果} & \bf{检测值} & \bf{参考范围} & \bf{功能}\\
\hline
\parbox[c]{\hsize}{\vskip10pt\centering 罗斯拜瑞氏菌属(Roseburia) \vskip10pt} &
正常 &
8.37
& ≥ 2.32
& 产生丁酸等有益物质，抑制肠道炎症，有利于肠道及人体健康 \\
\hline
\parbox[c]{\hsize}{\vskip10pt\centering 柔嫩梭菌属(Faecalibacterium) \vskip10pt} &
正常 &
166.68
& ≥ 10.94
& 发酵纤维素产生丁酸等有益物质，抑制肠道炎症，促进肠道健康 \\
\hline
\parbox[c]{\hsize}{\vskip10pt\centering 双歧杆菌属(Bifidobacterium) \vskip10pt} &
正常 &
0.06
& ≥ 0.01
& 有益菌，降解人体不能消化的多糖，产乳酸，调节免疫及肠道环境 \\
\hline
\parbox[c]{\hsize}{\vskip10pt\centering 乳酸杆菌属(Lactobacillus) \vskip10pt} &
异常 &
0.00
\raisebox{-0.3ex}{\includegraphics[scale=1]{redarrow_down.pdf}}
& ≥ 0.01
& 肠道益生菌，能够生成乳酸，抑制有害菌及炎症，调节肠道环境 \\
\hline
\parbox[c]{\hsize}{\vskip10pt\centering 阿克曼氏菌属(Akkermansia) \vskip10pt} &
异常 &
0.00
\raisebox{-0.3ex}{\includegraphics[scale=1]{redarrow_down.pdf}}
& ≥ 0.01
& 降解粘蛋白、调节免疫，有利于肠黏膜完整性，保持正常体重 \\
\hline
\end{longtable}

\vspace*{0mm}

\begin{spacing}{1.25}
\noindent\fontsize{9pt}{11pt}\selectfont {综合您的肠道主要有益菌检测结果，乳酸杆菌属、阿克曼氏菌属的含量异常，
不利于抑制肠道炎症、调节肠道环境、调节免疫等，
需引起注意。} \\

\end{spacing}

\vspace*{8mm}

\noindent{\fontsize{15pt}{11pt}\selectfont{ 四、主要有害菌检测结果}}

%\vspace*{-4.25mm}
\fontsize{9pt}{11pt}\selectfont
\arrayrulecolor{gray2}
\begin{longtable}{|m{4.8cm}<{\centering}|m{1.3cm}<{\centering}|m{1.3cm}<{\centering}|m{1.3cm}<{\centering}|m{4.9cm}<{\centering}|}
\hline
\parbox[c]{\hsize}{\vskip10pt\centering\bf {检测内容} \vskip10pt} & \bf{检测结果} & \bf{检测值} & \bf{参考范围} & \bf{功能}\\
\hline
\parbox[c]{\hsize}{\vskip10pt\centering 脱硫弧菌属(Desulfovibrio) \vskip10pt} &
正常 &
2.79
& ≤ 4.13
& 产生硫化氢，刺激肠道产生炎症反应，不利于肠道健康 \\
\hline
\parbox[c]{\hsize}{\vskip10pt\centering 多尔氏菌属(Dorea) \vskip10pt} &
正常 &
1.23
& ≤ 10.09
& 肠道的主要产气菌之一，与肠易激综合征等疾病相关 \\
\hline
\parbox[c]{\hsize}{\vskip10pt\centering 弯曲杆菌属(Campylobacter) \vskip10pt} &
正常 &
0.00
& ≤ 0.20
& 多数菌种为致病菌，可引起弯曲菌病，表现为严重腹泻或痢疾综合征 \\
\hline
\end{longtable}

\begin{spacing}{1.25}
\noindent\fontsize{9pt}{11pt}\selectfont {综合您的肠道主要有害菌检测结果，各有害菌含量均处于正常范围，有害物质产生少、
肠道炎症风险低，整体状况良好。请您保持良好的饮食及生活习惯，促进肠道菌群平衡。} \\

\end{spacing}

\end{document}
